% ex: ts=2 sw=2 sts=2 et filetype=tex
% SPDX-License-Identifier: CC-BY-SA-4.0
\documentclass[aspectratio=169]{beamer}

\usetheme{Boadilla}

\setbeamertemplate{navigation symbols}{} % To remove the navigation symbols from the bottom of all slides

\graphicspath{{../../../slides/linux/img/}{../styles/iteso/}}

% use the command "usepackage" instead of "usecolortheme" to use the path to styles
\usepackage{../styles/iteso/beamercolorthemeiteso}

\usepackage[utf8]{inputenc}
\usepackage[T1]{fontenc} %% https://tex.stackexchange.com/questions/664/why-should-i-use-usepackaget1fontenc
\usepackage{graphicx} % Allows including images
\usepackage{listings}
\usepackage{colortbl}

\title{Secure Shell} % The short title appears at the bottom of every slide,
                                               % the full title is only on the title page
\author{Miguel Bernal Marin} % Your name
\institute[ITESO]
{
 ITESO, Universidad \\
 Jesuita de Guadalajara \\
\medskip
\textit{miguel.bernal@iteso.mx}
}
\date{
    \today
} % Date, can be changed to a custom date

%New colors defined below
\definecolor{codeNumbers}{rgb}{0.5,0.5,0.5}
\definecolor{editorGray}{rgb}{0.95, 0.95, 0.95}
\definecolor{editorOcher}{rgb}{1, 0.5, 0} % #FF7F00 -> rgb(239, 169, 0)
\definecolor{editorGreen}{rgb}{0, 0.5, 0} % #007C00 -> rgb(0, 124, 0)
\definecolor{brown}{rgb}{0.69,0.31,0.31}

%Code listing style named "mystyle"
\lstdefinestyle{mystyle}{
  % General design
  backgroundcolor=\color{editorGray},
  basicstyle={\ttfamily},
  numberstyle=\tiny\color{codeNumbers},
  xleftmargin={0.2cm},
  numbersep=5pt,
  numbers=left,
  stepnumber=1,
  firstnumber=1,
  numberfirstline=true,
  % Code design
  identifierstyle=\color{black},
  keywordstyle=\color{blue}\bfseries,
  ndkeywordstyle=\color{editorGreen}\bfseries,
  stringstyle=\color{editorOcher}\ttfamily,
  commentstyle=\color{brown}\ttfamily,
  % Code
  alsodigit={.:;},
  tabsize=2,
  showtabs=false,
  showspaces=false,
  showstringspaces=false,
  extendedchars=true,
  breaklines=true,
  % Español
  literate=%
  {á}{{\'a}}1 {é}{{\'e}}1 {í}{{\'i}}1 {ó}{{\'o}}1 {ú}{{\'u}}1
  {Á}{{\'A}}1 {É}{{\'E}}1 {Í}{{\'I}}1 {Ó}{{\'O}}1 {Ú}{{\'U}}1
  {¡}{{!`}}1  {¿}{{?`}}1
  {ñ}{{\~n}}1 {Ñ}{{\~N}}1
}

%"mystyle" code listing set
\lstset{style=mystyle}

%------------------------------------------------------------
%The next block of commands puts the table of contents at the
%beginning of each section and highlights the current section:
\AtBeginSection[]
{
  \begin{frame}
    \frametitle{Contenido}
    \tableofcontents[currentsection]
  \end{frame}
}
%------------------------------------------------------------

\newcommand{\pausa}{\pause} % Para usar una pausa en las presentaciones
%\newcommand{\pausa}{}      % Para que NO salgan las pausas

\begin{document}
{ % Portada de la universidad
  \setbeamertemplate{footline}{}
  \usebackgroundtemplate{\includegraphics[width=\paperwidth]{fondo-portada.png}}
  \begin{frame}
  \end{frame}
}

\begin{frame}
    \titlepage
\end{frame}

\usebackgroundtemplate{\includegraphics[width=\paperwidth]{fondo.png}}

\begin{frame}
    \frametitle{Contenido}
    \tableofcontents
\end{frame}

% ex: ts=2 sw=2 sts=2 et filetype=tex
% SPDX-License-Identifier: CC-BY-SA-4.0

\section{Introducción a Wireguard}

\begin{frame}[c]{¿Qué es Wireguard?}
    \begin{columns}
        \column{0.5\textwidth}
        WireGuard es una aplicación de software libre y de código abierto
        y un protocolo de comunicación que implementa técnicas de red
        privada virtual (\textbf{VPN}) para crear \emph{conexiones seguras
        punto a punto} en configuraciones enrutadas o puenteadas.
        \column{0.5\textwidth}
        \begin{center}
            \includegraphics[scale=0.5]{wireguard-logo.png}
        \end{center}
    \end{columns}
\end{frame}

\begin{frame}[c]{¿Qué es Wireguard? (continuación)}
    \begin{columns}
        \column{0.5\textwidth}
        \begin{center}
            \includegraphics[scale=0.5]{kr-donenfeld.jpg}
        \end{center}
        \column{0.5\textwidth}
        Se ejecuta como un módulo dentro del núcleo \textbf{Linux} y tiene
        como objetivo un mejor rendimiento que los protocolos de tunelización
        IPsec y OpenVPN.

        \vspace{\baselineskip}
        Fue escrito por \textbf{Jason A. Donenfeld} y se publica bajo la
        segunda versión de la GNU General Public License (GPL).
    \end{columns}
\end{frame}

\begin{frame}[c]{Características}
  \begin{itemize}
    \item Tiene como objetivo proporcionar una VPN que sea
          simple y altamente efectiva.
    \pausa
    \item Código fuente de alrededor de 4000 líneas.
    \pausa
    \item Utiliza criptografía de última generación, como el \emph{framework}
          del protocolo Noise, Curve25519, ChaCha20, Poly1305, BLAKE2,
          SipHash24, HKDF y construcciones seguras y confiables.

  \end{itemize}
\end{frame}

\section{Creación de llaves}

\begin{frame}[fragile]
  \frametitle{Generando llave privada}

  WireGuard requiere claves públicas y privadas codificadas en base64.
  Estos se pueden generar utilizando la utilidad \textbf{wg(8)}: 

  \vspace{\baselineskip}
  \begin{lstlisting}[language=Bash]
  $ umask 077
  $ wg genkey > privatekey
  \end{lstlisting}

  \vspace{\baselineskip}
  esto creará una clave privada en la salida estándar que contiene una
  nueva clave privada. 
\end{frame}

\begin{frame}[fragile]
  \frametitle{Generando llave publica}

  Luego puede derivar su clave pública de su clave privada:

  \vspace{\baselineskip}
  \begin{lstlisting}[language=bash]
  $ wg pubkey < privatekey > publickey
  $
  \end{lstlisting}

  \vspace{\baselineskip}
  Esto leerá la clave privada de la entrada estándar y escribirá la
  clave pública correspondiente en la clave pública en la salida estándar.
\end{frame}

\begin{frame}[fragile]
  \frametitle{Generando ambas llaves a la vez}

  Por supuesto, puedes hacer todo esto a la vez: 

  \vspace{\baselineskip}
  \begin{lstlisting}[language=Bash]
  $ umask 077
  $ wg genkey | tee privatekey | wg pubkey > publickey
  \end{lstlisting}

\end{frame}

\section{Configurando las interfaces de red}

\begin{frame}[fragile]
  \frametitle{Agregando una interfaz}

  Se puede agregar una nueva interfaz a través de \textbf{ip-link(8)},
  que debería manejar automáticamente la carga del módulo:

  \vspace{\baselineskip}
  \begin{lstlisting}[language=Bash]
  # ip link add dev wg0 type wireguard
  \end{lstlisting}

  Los usuarios que no son de Linux escribirán en su lugar wireguard-go wg0.
\end{frame}

\begin{frame}[fragile]
  \frametitle{Asignando IP}

  Se puede asignar una dirección IP y un "peer" con \textbf{ip-address(8)}

  \vspace{\baselineskip}
  \begin{lstlisting}[language=Bash]
  # ip address add dev wg0 192.168.2.1/24
  \end{lstlisting}

  \vspace{\baselineskip}
  O, si solo hay dos pares en total, algo como esto podría ser más deseable:

  \vspace{\baselineskip}
  \begin{lstlisting}[language=Bash]
  # ip address add dev wg0 192.168.2.1 peer 192.168.2.2
  \end{lstlisting}
\end{frame}

\begin{frame}[fragile]
  \frametitle{}

  La interfaz se puede configurar con claves y puntos finales del
  mismo nivel con la utilidad \textbf{wg(8)} incluida: 

  \vspace{\baselineskip}
  \begin{lstlisting}[language=Bash]
  # wg setconf wg0 myconfig.conf
  \end{lstlisting}

  \vspace{\baselineskip}
  o con

  \vspace{\baselineskip}
  \begin{lstlisting}[language=Bash]
  # wg set wg0 listen-port 51820 private-key /path/to/private-key peer ABCDEF... allowed-ips 192.168.88.0/24 endpoint 209.202.254.14:8172
  \end{lstlisting}
\end{frame}

\begin{frame}[fragile]
  \frametitle{}

  Finalmente, la interfaz se puede activar con  \textbf{ip-link(8)}: 

  \vspace{\baselineskip}
  \begin{lstlisting}[language=Bash]
  # ip link set up dev wg0
  \end{lstlisting}

  \vspace{\baselineskip}
  También existen los comandos \textbf{wg show} y \textbf{wg showconf}, para ver
  la configuración actual. Llamar a wg sin argumentos por defecto es 
  lamar a wg show en todas las interfaces de WireGuard.

\end{frame}

%\begin{frame}[c]{Caracteres de escape}
%\end{frame}

%\begin{frame}[fragile]
%  \frametitle{Cadenas de texto}
%
%  \vspace{\baselineskip}
%  \textcolor{codeKeyword}{print}():
%
%  \vspace{\baselineskip}
%  \begin{lstlisting}[language=Python]
%  \end{lstlisting}
%\end{frame}


{ % Portada de la universidad
  \setbeamertemplate{footline}{}
  \usebackgroundtemplate{\includegraphics[width=\paperwidth]{fondo-portada.png}}
  \begin{frame}
  \end{frame}
}
\end{document}
